\section{Double Chocolate Chip Muffins}

\begin{quotex}
Venus is the source of the power of taste. She governs the stomach whence emerges the ability and taste for eating and drinking.
\flright{\textit{Picatrix}}
\end{quotex}

On Venus day, here is an indulgence for you. For several years, I worked in the laboratory of my father's chemical business. There I learned special skills that led to my newfound affinity for baking muffins. I had to determine the percentage of ash in phenolic resins by baking a sample of resin and then measuring what was left. I had to mix and I had to measure viscosity. In the end I had to clean the glassware by soaking them in sulfuric acid. I've been informed that doing so is not recommended in the kitchen, so I make sure there is someone around to clean up.

During the lockdown, I was limited and could not visit my favorite bakeries. Hence, I took up baking. Although I used the metric system in the laboratory, this is not suitable for the American kitchen. All recipes are measured in traditional units as are kitchen supplies. So in the following recipe, I provide the American measures. I looked up the metric equivalents which you can try, but I suggest verifying the conversions for yourself. Although lab work is quite precise, the culinary arts allow some leeway for poetic flair.

I suppose I can explain the chemistry of baking or why baking powder is used or the purpose of the eggs, but it won't affect the outcome. Happy Venus Day!

\paragraph{Ingredients}

Metric equivalents are within parentheses.

\begin{itemize}


\item 2 cups (320g) all-purpose flour

\item 1 cup (200g) sugar

\item 1 Tablespoon (15ml) baking powder

\item 1/2 cup (56g) unsweetened cocoa

\item A little salt

\item 1/2 cup (112g) melted butter (8 tbsp or one stick)

\item 2 large eggs, beaten (a US large egg is 57 g. A European large egg is 63 g, so close enough)

\item 3/4 cup (180 ml) milk

\item A little Vanilla (optional)

\item 1 to 2 cups by volume (250 to 500ml) of chocolate chips
\end{itemize}


I suggest 1 cup of chocolate chips for a more understated muffin. The package directions recommend 2 cups, but they just want you to use up the chips as fast as you can.

\paragraph{Mixing}


Combine all the dry ingredients in a mixing bowl and combine thoroughly.

In the meantime, melt the butter, either stovetop or microwave (preferred), and beat the eggs.

Add the wet ingredients (milk, butter, eggs) and mix thoroughly. It should turn into a pleasantly colored dark brown, viscous paste. Now add the chocolate chips and mix until they are randomly distributed.

\paragraph{Baking}


Spray some oil into a muffin pan and put some of the muffin dough into each of the holes. An ice cream scoop works well, although I now use two spoons: one to pick up some dough and the other to scrape it into the howl

Bake at 375°F (190°C) about 22 (22) minutes I use a Ninja toaster oven which heats up quite quickly and holds heat well. You may need to cook longer. A trick that looks really cool if you are trying to impress friends is to insert a wooden toothpick into the muffin. If it comes out clean, the muffin is ready; otherwise, it needs longer. In practice, the technique is not very helpful, so I just look see if the top is encrusted.

This makes a dozen (12) American muffins or 10 metric muffins.


\flrightit{Posted on 2021-09-17 by Cologero}

\begin{center}* * *\end{center}

\begin{footnotesize}
\begin{sffamily}
\texttt{Max Leyf on 2021-09-17 at 23:23 said:}

It's an alchemical allegory

\hfill

\texttt{Ismo Meinander on 2021-09-18 at 13:36 said:}

Max Leyf: Cooking is a fine allegory of alchemy in general. From the raw ingredients you can make gold if you mix, separate and unite (solve \& coagula) them the right way.

What would a Venusian diet look like as a whole? Bread and wine?

\hfill

\texttt{Tannheuser on 2021-09-21 at 11:23 said:}

The dry ingredients must first be isolated, while the wet ingredients -- milk, butter and eggs must be be mixed with the chocolate chips ? and cooked until they are suffused with the fiery nature of chocolate ?. The next phase cannot be revealed here, but those who are able will discover the rest on their own.

\hfill

\texttt{Mike on 2021-09-21 at 21:46 said:}

You'll have to taste it for yourself


\hfill

\end{sffamily}
\end{footnotesize}

